% ============================================================================
% CHAPITRE 3 — SECTION 3.1 : INTRODUCTION
% Fichier : partie3/introduction_validation.tex
% ============================================================================

\section{INTRODUCTION}
\label{sec:introduction_experimentation}
% ============================================================================
\begin{sloppypar}

Les Trilemmes CRO classique et post-quantique \'etablis au
chapitre~\ref{chap:cadre_formel} sont des r\'esultats conditionnels, car leur
port\'ee pratique d\'epend enti\`erement de la possibilit\'e de mesurer
empiriquement les param\`etres institutionnels qui les param\`etrent. Pour le
cadre classique (Contribution~C2), il s'agit de $\Delta$ (capacit\'e
d'extraction), $\alpha$ (granularit\'e des verdicts), $\beta$ (r\'egularit\'e)
et $H(\mathcal{C})$ (entropie contextuelle). Pour l'extension post-quantique
(Contribution~C3), s'y ajoutent $\Delta_q$ (capacit\'e quantique) et $\eta_q$
(coefficient de perte d'interpr\'etabilit\'e quantique). Sans proc\'edure de
mesure, ces grandeurs demeurent des postulats~; sans estimation, les bornes
demeurent muettes sur les syst\`emes r\'eels.

\subsection{Statut \'epist\'emique des deux volets}
\label{subsec:statut_epistemique}

Ce chapitre r\'epond \`a cette double exigence en deux temps de nature
distincte, qu'il importe de ne pas confondre.

\textbf{Premi\`erement}, les
sections~\ref{sec:protocoles_experimentaux}
\`a~\ref{subsec:calendrier_ressources} sp\'ecifient, jusqu'au niveau
op\'erationnel, le dispositif exp\'erimental permettant de mesurer chacun des
param\`etres, qu'il s'agisse de la stratification des \'echantillons, des estimateurs, des crit\`eres de
fiabilit\'e inter-\'evaluateurs ou des seuils de falsification. Ce programme n'a pas
\'et\'e d\'eploy\'e dans le cadre de cette th\`ese~; son statut exact est
explicit\'e \`a la sous-section~\ref{subsubsec:statut_deploiement}, qui referme
cette partie.

\textbf{Deuxi\`emement}, les
sections~\ref{sec:derivation_opposabilite}
\`a~\ref{sec:reproductibilite} proc\`edent \`a une analyse
\emph{r\'etrospective et analytique}. Elle d\'erive d'abord, pour trois
syst\`emes cryptographiques r\'eellement d\'eploy\'es, les valeurs
d'opposabilit\'e et les produits CRO qui servent de r\'ef\'erence \`a tout le
chapitre. Elle estime ensuite, par une voie enti\`erement distincte (triangulation de la litt\'erature cognitive, de l'analyse documentaire et des
sp\'ecifications techniques), les param\`etres institutionnels qui
param\`etrent les bornes avant de mettre les deux estimations en regard. Ces deux voies
ne sont pas ind\'ependantes, puisque les valeurs de $\Delta$ ont d'abord \'et\'e fix\'ees
par inversion de la borne. Les syst\`emes d\'eploy\'es ne sont pas non plus des
instances du th\'eor\`eme (Remarque~\ref{rem:portee_cas_etude}). Cette partie
constitue donc une \emph{calibration illustrative}. Elle v\'erifie que les ordres
de grandeur retenus sont mutuellement coh\'erents, sans constituer une validation
empirique.

Le chapitre sp\'ecifie d'abord les protocoles de mesure puis exploite les donn\'ees existantes.
Cet ordre m\'ethodologique r\'epond \`a une n\'ecessit\'e logique, car on ne peut extraire
d'information de mesures que si l'on sait d\'efinir ce qui doit \^etre mesur\'e.
La sp\'ecification des protocoles pr\'ec\`ede donc l'application aux
syst\`emes d\'eploy\'es. Les estimations r\'etrospectives restent ainsi ancr\'ees
dans une proc\'edure explicite et falsifiable.

\subsection{Param\`etres \`a mesurer et statut de leur validation}
\label{subsec:parametres_mesurer}

Le tableau~\ref{tab:parametres_mesure} r\'ecapitule, pour chaque contribution,
les param\`etres concern\'es, le protocole de mesure qui leur est d\'edi\'e et
l'\'etat de leur validation.

\begin{table}[htbp]
\centering
\caption{Param\`etres institutionnels~: protocole de mesure et statut de validation}
\label{tab:parametres_mesure}
\small
\begin{tabular}{@{}lllll@{}}
\toprule
\textbf{Contribution} & \textbf{Param\`etres} & \textbf{Protocole de mesure} &
\textbf{Statut de validation} & \textbf{Type} \\
\midrule
C2 (Trilemme classique)
  & $\Delta$, $\alpha$, $\beta$, $H(\mathcal{C})$
  & \S\ref{subsec:protocole_delta}--\S\ref{subsec:protocole_conjoint}
  & Calibration illustrative & R\'etrospectif \\
C3 (Trilemme post-quantique)
  & $\Delta_q$, $\eta_q$, $\beta_q$
  & \S\ref{subsec:protocole_delta_q}--\S\ref{subsec:protocole_combine}
  & Calibration illustrative & R\'etrospectif \\
C4 (Famille SOTP)
  & Opposabilit\'e r\'ealis\'ee
  & \S\ref{subsec:protocole_opp}
  & Calibration illustrative & R\'etrospectif \\
\bottomrule
\end{tabular}
\end{table}

La section~\ref{subsec:protocole_delta} d\'etaille la mesure de $\Delta$,
la section~\ref{subsec:protocole_HC} celle de $H(\mathcal{C})$ tandis que la
section~\ref{subsec:protocole_opp} porte sur la mesure de l'opposabilit\'e sur
syst\`emes d\'eploy\'es. Les protocoles post-quantiques
($\Delta_q$ en section~\ref{subsec:protocole_delta_q}, $\eta_q$ en
section~\ref{subsec:protocole_eta_q}) compl\`etent le dispositif, avant que
la section~\ref{subsec:protocole_combine} n'en articule la validation
conjointe.

La section~\ref{sec:derivation_opposabilite} \'etablit ensuite les
opposabilit\'es des trois syst\`emes classiques. La
section~\ref{sec:validation_retrospective} met ces d\'erivations en regard
des param\`etres institutionnels estim\'es par ailleurs. Cette mise en regard
v\'erifie la coh\'erence des ordres de grandeur. Le test du cadre rel\`eve du
programme exp\'erimental (crit\`eres de r\'efutation du
tableau~\ref{tab:criteres_refutation}).

\subsection{Table unique des param\`etres de r\'ef\'erence}
\label{subsec:parametres_reference}

Le tableau~\ref{tab:parametres_reference} fixe la seule source de r\'ef\'erence
pour $H(\mathcal{C})$, $\alpha$, $\Delta$ et $\Delta_q$. Toute valeur employ\'ee
ailleurs y renvoie. Il est reproduit dans le fichier
\texttt{06\_calcul/data/parametres\_reference.csv} du d\'ep\^ot
compagnon~\cite{CROTrilemmaDepot}.

\begin{table}[htbp]
\centering
\caption{Param\`etres institutionnels de r\'ef\'erence (valeurs illustratives)}
\label{tab:parametres_reference}
\small
\begin{tabular}{@{}lccccc@{}}
\toprule
\textbf{Syst\`eme} & $H(\mathcal{C})$ domaine & $H(\mathcal{C})$ esp\`ece & $\alpha$ & $\Delta$ & $\Delta_q$ / $\eta_q$ \\
\midrule
ECDSA (contrat) & 20 & 20 & 1{,}0 & 35 & --- \\
Helios (scrutin) & 40 & 8 & 1{,}2 & 56 & --- \\
zk-SNARKs (AML) & 70 & 11 & 1{,}5 & 73 & --- \\
ML-DSA-65 (contrat) & 20 & 10{,}5 & 1{,}0 & --- & 28 (hyp.) / $0$ \\
\bottomrule
\end{tabular}

\smallskip
\begin{minipage}{0.92\linewidth}\footnotesize
Entropies en bits. \og Domaine\fg{} d\'esigne la richesse du contexte institutionnel
g\'en\'eral (chapitre~\ref{chap:cadre_formel})~; \og esp\`ece\fg{} d\'esigne le contexte d'une
d\'ecision particuli\`ere (section~\ref{sec:derivation_opposabilite}). Pour
ML-DSA-65, $\eta_q = 0$ car la signature est une cha\^{\i}ne de bits classique
(les valeurs $\eta_q > 0$ du chapitre d\'ecrivent un sc\'enario hypoth\'etique)~;
$\Delta_q = 28$~bits est une hypoth\`ese de travail
(section~\ref{sec:etude_cas_dilithium}), \`a estimer par le protocole de la
section~\ref{subsec:protocole_delta_q}.
\end{minipage}
\end{table}

\end{sloppypar}